\subsection{管与横截面}

引理 4. 设 H 是紧李群，\(\rho : \mathrm{H} \to \mathbf{GL}(\mathrm{V})\) 是 H 在有限维实向量空间 V 上的连续（因而解析）表示，W 是 V 中原点的一个邻域。存在原点的一个开邻域 B，它包含于 W 且在 H 的作用下稳定；还存在一个解析同构 \(u: \mathrm{V} \to \mathrm{B}\)，与 H 的作用相容，并满足 \(u(0) = 0\) 和 \(\mathrm{Du}(0) = \mathrm{Id}_{\mathrm{V}}\)。

在 V 上选取一个 H-不变内积（§1, no. 1）。存在实数 \(r > 0\)，使得半径为 \(r\) 的开球 B 包含于 W 中；它显然在 H 的作用下稳定。置 \(u(v) = r(r^2 + \|v\|^2)^{-1/2} v\)，对所有 \(v \in V\) 成立；于是 \(u\) 是从 V 到 B 的双射解析映射，与 H 的作用相容，且其逆映射 \(w \mapsto r(r^2 - \|w\|^2)^{-1/2} w\) 是解析的。此外，\(u(0) = 0\) 且 \(\mathrm{D}u(0) = \mathrm{Id}_{\mathrm{V}}\)。

命题 5. 设 H 是紧李群，\((h,x)\mapsto hx\) 是 H 在 X 上的 \(C^{r}\) 类左作用律，x 是 X 中在 H 的作用下固定的一点。则群 H 在线性地作用于向量空间 \(T=T_{x}(X)\) 上；存在一个开嵌入 \(\varphi:T\to X\)，为 \(C^{r}\) 类，与 H 的作用相容，并满足 \(\varphi(0)=x\) 且 \(T_{0}(\varphi)\) 是 T 的恒等映射。

设 \((\mathrm{U},\psi,\mathrm{E})\) 是 X 在 x 处的一个图册，使得 U 在 H 的作用下稳定（no. 2, Lemma 3）且 \(\psi(x)=0\)。通过 \(\mathrm{T}_{x}(\psi)\) 将 E 与 T 识别，并置

\[
\psi^ {\sharp} (y) = \int_ {\mathrm{H}} h. \psi (h ^ {- 1} y) d h \quad \text { for } y \in \mathrm{U},
\]

其中 dh 是 H 上总质量为 1 的 Haar 测度。

于是（§6, no. 4, Cor. 1）\(\psi^{\sharp}\) 是从 U 到 T 的 \(\mathrm{C}^r\) 类态射，与 H 的作用相容，并满足 \(\psi^{\sharp}(x) = 0\) 和 \(d_x\psi^\sharp = \mathrm{Id}_\mathrm{T}\)。因此存在一个开集 \(\mathrm{U}' \subset \mathrm{U}\)，包含 \(x\)，以及 T 中 0 的一个开邻域 V，使得 \(\psi^{\sharp}\) 诱导出同构 \(\theta: \mathrm{U}' \to \mathrm{V}\)。必要时限制 \(\mathrm{U}'\) 和 V，可假设它们在 H 的作用下稳定，并且存在与 H 的作用相容的同构 \(u: \mathrm{T} \to \mathrm{V}\)（引理 4）。现在只需取 \(\varphi = \theta^{-1} \circ u\)。

回忆（Differentiable and Analytic Manifolds, Results, 6.5.1），若 G 是李群，H 是 G 的李子群，且 Y 是 H 在其上左作用的流形，则记 \(G \times^{H}Y\) 为乘积流形 \(G \times Y\) 关于 H 的右作用 \(((g,y),h) \mapsto (gh,h^{-1}y)\) 的商；这是一个李群 G 在其上自然地左作用的流形；投影 \(G \times^{H}Y \to G/H\) 是纤维为 Y 的丛。此外，若 Y 是 H 在线性作用的有限维向量空间，则 \(G \times^{H}Y\) 自然具有以 G/H 为底空间的向量 G-丛结构（Differentiable and Analytic Manifolds, Results, 7.10.2）。

设 G 是一个在流形 X 上真作用的李群（General Topology, Chap. III, §4, no. 1, Def. 1），且作用律 \((g, x) \mapsto gx\) 是 \(C^{r}\) 类的。则对于 X 中的每一点 x，x 的轨道 Gx 是 X 的闭子流形，同构于李齐性空间 \(G/G_{x}\)，其中 \(G_{x}\) 是 G 中 x 的固定子群（cf.~Chap. III, §1, no. 7, Prop. 14 (ii), and General Topology, Chap. III, §4, no. 2, Prop. 4）；后者是紧李群（loc. cit.）。

命题 6. 假设流形 X 是仿紧的；设 x 是 X 中的一点，\(G_{x}\) 是其固定子群。存在有限维解析线性表示 \(\tau: G_{x} \to \mathbf{GL}(W)\)，以及一个开嵌入 \(\alpha: G \times ^{G_{x}} W \to X\)，为 \(C^{r}\) 类，与 G 的作用相容，并将 \((e,0) \in G \times W\) 的等价类映到 x。

置 \(T = T_{x}(X)\)。设 W 是 T 中 \(T_{x}(Gx)\) 的一个补子空间，在 \(G_{x}\) 的作用下稳定（例如，取 \(T_{x}(Gx)\) 相对于 T 上 \(G_{x}\)-不变内积的正交补）。另一方面，设 \(\varphi : T \to X\) 是具有命题 5 所述性质的态射（相对于 \(H = G_{x}\)）。考虑态射 \(\lambda : G \times W \to X\)，其定义为 \(\lambda(g, w) = g\varphi(w)\)。它通过取商诱导出态射 \(\mu : G \times^{G_{x}}W \to X\)，为 \(C^{r}\) 类，与 G 的作用相容，并将 \((e, 0)\) 的等价类 z 映到 x。

我们证明 \(\mu\) 在点 \(z\) 处是局部微分同胚的。有

\[
\begin{array}{c} \dim (\mathrm{G} \times {} ^ {\mathrm{G} _ {x}} \mathrm{W}) = \dim (\mathrm{G}) + \dim (\mathrm{W}) - \dim (\mathrm{G} _ {x}) \\ = \dim (\mathrm{G} x) + \dim (\mathrm{W}) = \dim (\mathrm{T}), \end{array}
\]

所以只需证明 \(\mu\) 在 \(z\) 处是下沉映射，或等价地，\(\lambda\) 在 \((e,0)\) 处是下沉映射。但是，切映射 \(\mathrm{T}_{(e,0)}(\lambda):\mathrm{T}_e(\mathrm{G})\oplus \mathrm{W}\to \mathrm{T}\) 等于 \(\mathrm{T}_e(\rho (x)) + i\)，其中 \(\rho (x)\) 是轨道映射 \(g\mapsto gx\)，而 \(i\) 是从 W 到 T 的典则嵌入；由于 \(\operatorname {Im}\operatorname {T}_e(\rho (x)) = \operatorname {T}_x(\operatorname {G}x)\)，切映射 \(\mathrm{T}_{(e,0)}(\lambda)\) 是满射，且 \(\mu\) 在 \(z\) 处是局部微分同胚的。

我们将证明，存在一个开邻域 \(\Omega\)，它是 \(Gz\) 在 \(G \times ^{G_x}W\) 中的邻域，在 \(G\) 的作用下稳定，使得 \(\mu\) 将 \(\Omega\) 同构到 \(X\) 的一个开子集上。这将推出该命题：事实上，\(\Omega\) 在 \(G \times W\) 中的逆像在 \(G\) 的作用下稳定，因而形如 \(G \times B\)，其中 \(B\) 是 \(W\) 中包含原点且在 \(G_x\) 的作用下稳定的开子集；必要时限制 \(\Omega\)，可假设存在同构 \(u: W \to B\)，与 \(G_x\) 的作用相容（引理 4）。显然，复合态射 \(\alpha: G \times ^{G_x}W \xrightarrow{(\text{Id}, u)} G \times ^{G_x}B \xrightarrow{\mu} X\) 满足命题中的条件。

因此，该命题是下列引理的推论：

引理 5. 设 Z 是一个 \(C^{r}\) 类分离流形，带有一个左作用律 \(m: G \times Z \to Z\)，为 \(C^{r}\) 类，且 \(\mu: Z \to X\) 是一个与 G 的作用相容的（\(C^{r}\) 类）态射。设 z 是 Z 中的一点，且 \(x = \mu(z)\)。假设 \(\mu\) 在 z 处是局部微分同胚的，且 G 中 z 的固定子群等于 x 的固定子群 \(G_{x}\)。则存在一个开邻域 \(\Omega\)，它是轨道 Gz 的邻域，在 G 的作用下稳定，使得 \(\mu\) 将 \(\Omega\) 同构到 X 的一个开子集上。

由于 \(\mu\) 与 G 的作用相容，它在 Gz 的每一点处都是局部微分同胚的；由于典则映射 \(G/G_{x} \rightarrow Gx\) 是同胚，由 \(\mu\) 诱导的从 Gz 到 Gx 的映射也是同胚。因此由 no. 1 的命题 2 可知，Z 中存在 Gz 的一个开邻域 U，使得 \(\mu\) 将 U 开嵌入 X。

由于 G 在 X 上真作用，存在 x 的一个开邻域 V 以及 G 的一个紧子集 K，使得当 \(gV \cap V = \varnothing\)（\(g \notin K\) 时）；特别地，\(e \in K\)。集合 \(W_{1}\) 由 Z 中满足 \(y \in Z\) 且 \(Ky \subset U\) 的点组成，是开的：事实上，\(Z - W_{1}\) 是闭集 \((\mathrm{K} \times \mathrm{Z}) - m^{-1}(\mathrm{U})\) 在真投影 \(pr_{2}: K \times Z \to Z\) 下的像。置 \(W = W_{1} \cap \mu^{-1}(V)\)；这是 Z 的一个包含 z 的开子集，并满足下列条件：

\begin{enumerate}
\def\labelenumi{(\roman{enumi})}
\item
  KW \(\subset\) U，特别是 W \(\subset\) U；
\item
  \(\mu(\mathrm{W}) \subset \mathrm{V}.\)
\end{enumerate}

置 \(\Omega = \mathrm{GW}\)，并考虑 \(\mu\) 在 \(\Omega\) 上的限制。这是一个局部微分同胚态射，因为 \(\Omega\) 的每一点都在 \(\mathbf{G}\) 的作用下与 \(\mathbf{U}\) 中的一点共轭。我们证明它是单射：设 \(g, h\) 属于 \(\mathbf{G}\)，\(u, v\) 属于 \(\mathbf{W}\)，且满足 \(\mu(gu) = \mu(hv)\)。置 \(k = g^{-1}h\)；则 \(\mu(u) = k\mu(v)\)，由 (ii) 得 \(k \in \mathbf{K}\)。但由 (i)，\(kv\) 和 \(u\) 都属于 \(\mathbf{U}\)；于是 \(u = kv\)，因为 \(\mu\) 在 \(\mathbf{V}\) 上的限制是单射，所以 \(gu = hv\)。因此 \(\mu\) 在 \(\Omega\) 上的限制是单射，从而根据 Differentiable and Analytic Manifolds, Results, 5.7.8，它是到 \(\mathbf{X}\) 的一个开子流形上的同构，证明完成。

在命题 6 的条件下，\(\alpha\) 的像是轨道 A 的点 \(x\) 的一个开邻域 T，它带有以 A 为底空间的向量丛结构，且零截面就是轨道 A 本身。对于每个点 \(a \in \mathrm{A}\)，该向量丛的纤维 \(\mathrm{Y}_a\) 是 X 的子流形，在固定子群 \(\mathrm{G}_a\)（对应点 \(a\)）的作用下稳定，并且有一个从 \(\mathrm{G} \times^{\mathrm{G}_a}\mathrm{Y}_a\) 到 X 的态射，它将 \((g,y) \in \mathrm{G} \times \mathrm{Y}_a\) 的等价类映到 \(gy \in \mathrm{X}\)，并且它诱导出一个 \(\mathrm{C}^r\) 类态射，从 \(\mathrm{G} \times^{\mathrm{G}_a}\mathrm{Y}_a\) 到 T。称 \(\mathrm{Y}_a\) 为管 T 在 \(a\) 处的横截面。注意，在 \(a\) 处，\(\mathrm{Y}_a\) 的切空间典则同构于 \(\mathrm{Y}_a\)，并且它是 \(\mathrm{T}_a(\mathrm{A})\) 在 \(\mathrm{T}_a(\mathrm{X})\) 中的一个补空间；因此，以 A 为底空间的向量丛 T 典则同构于 A 在 X 中的法丛（Differentiable and Analytic Manifolds, Results, 8.1.3）。

